\chapter{Limites et menaces sur la validité}
\label{chap:10-limites}

Ce chapitre rassemble les limites du travail, pour qu'aucun lecteur n'ait à les chercher. Il les
range selon les quatre familles de menaces de \textcite{shadish2002experimental}. La validité
interne demande si l'effet mesuré vient bien de la cause supposée. La validité externe demande à
quelles populations, tâches et conditions le résultat s'étend. La validité de construit demande si
la mesure saisit la notion qu'elle prétend saisir. La validité de la conclusion statistique demande
si les tests et les intervalles autorisent la conclusion tirée. Une dernière section dit ce que
testerait un nouvel entraînement complet sur une autre machine.

\section{Validité interne}
\label{sec:val-interne}

\subsection{Les lecteurs du sondage sont des modèles d'IA agentique}

Tous les articles ont été lus par des agents IA, copies d'un même modèle
(\texttt{claude-opus-5-5}). Deux lecteurs ont fini sur un modèle voisin (\texttt{claude-opus-5})
pour treize articles, après l'arrêt d'une réponse par un filtre de contenu du fournisseur (écart 28 du sondage).
Aucun humain n'a relu les articles à l'aveugle : la relecture prévue par l'écart 13 a été abandonnée
avant de commencer (écart 23). Cinq garde-fous réduisent le risque sans l'annuler.

\begin{enumerate}
  \item \textbf{La citation obligatoire.} Chaque décision \enquote{oui} s'appuie sur une citation
        exacte et sa page. Chacun peut vérifier chaque décision contre le PDF. Le script
        \file{study/check\_codes.py} vérifie mécaniquement que chaque citation figure dans le texte.
  \item \textbf{L'audit des citations.} Après les premiers pourcentages, nous avons cherché les
        citations que la vérification mécanique ne retrouvait pas (écart 15, fichier
        \file{study/v1\_quotes\_audit.md}) : sur 89 citations, aucune n'est inventée, 83 figurent mot pour mot dans
        l'article, une est une paraphrase fidèle, 5 donnent des valeurs contraires aux règles.
  \item \textbf{Les doubles lectures à l'aveugle.} Vingt articles du sondage ont été relus par un
        second agent qui ignorait la première lecture. Sur les dix articles relus du premier tirage,
        les deux lecteurs s'accordent sur 127 décisions sur 130 (98~\%). Deux agents s'accordent sur
        au moins 80~\% des articles pour chaque critère, à une exception près, la mesure de
        performance (Y1), à 60~\% sur le premier tirage, avant que ses règles soient précisées. Dix
        articles de la trajectoire annuelle, relus de même, concordent à 80~\% au moins sur chacun
        des 23 critères (\cref{tab:accord}).
  \item \textbf{La recherche contre les \enquote{non}.} Un \enquote{non} ne s'appuie sur aucune
        citation. Chaque \enquote{non} des critères des verdicts a été confronté au texte de
        l'article avec des mots-clés fixés d'avance, puis chaque extrait trouvé a été jugé par un
        arbitre qui ne voyait ni l'article ni l'année (écart 27). Sur 377 décisions \enquote{non},
        10 passent à \enquote{oui} (2 en 2018, 8 en 2026) et aucun verdict ne change. Si l'on
        retourne aussi les 24 décisions jugées douteuses, l'incertitude (X2) et la sensibilité (Y9)
        passent juste au-dessus du seuil ($p = 0{,}07$ et $p = 0{,}054$).
  \item \textbf{Le retrait du modèle voisin.} Sans les treize articles lus en partie par
        \texttt{claude-opus-5}, les verdicts tiennent, sauf l'incertitude rapportée (X2), qui ne
        passe plus la correction de Holm ($p = 0{,}07$, écart 29).
\end{enumerate}

Ces garde-fous laissent une menace entière. Les deux lecteurs d'une double lecture sont des copies
du même modèle : leur accord mesure la stabilité de ce modèle, non l'accord de deux lecteurs
indépendants. Un biais systématique du modèle, partagé par toutes ses copies, échappe aux doubles
lectures. La consigne exacte, le modèle et les dates de chaque lecteur sont publiés
(\file{study/readers/}), pour qu'un autre lecteur, humain ou non, refasse la lecture.

\subsection{Les décisions prises après les résultats}

Plusieurs règles ont été fixées après avoir vu des résultats : la règle de verdict du sondage
(écart 20), l'harmonisation de la calibration de 2018 (écart 15), le classement descriptif sous
physique modifiée (écart 11 des simulations), les détails du bootstrap GRPO (écart 13) et la remise à
zéro du compte des échecs GRPO (écart 10). Le \cref{chap:09-ecarts} discute chacune. L'ordre
d'écriture des écarts antérieurs au 15 novembre 2026 repose sur notre parole et sur les dates des
commits (\cref{sec:parole}).

\subsection{Les essais des simulations}

Le registre reconstruit compte 42 tentatives arrêtées que le protocole ne déclarait pas
(\cref{sec:registres}). Chaque arrêt a frappé tous les essais en cours à la fois, sans lien avec un
score. Pour SAC et CrossQ, les courbes des tentatives arrêtées ont été écrasées : nous ne pouvons pas
montrer que la relance retrouve la même courbe. Le protocole annonçait Stable-Baselines3 2.6.0.
Tous les entraînements ont tourné avec la 2.9.2, la même pour tous les essais. Ces écarts ne
favorisent aucun programme. Ils sont dits ici pour que le lecteur les connaisse.

\subsection{Le matériel et la graine aléatoire du volet GRPO}

Pour Qwen, le matériel a changé avec la graine aléatoire : les deux essais à vraie récompense qui
donnent le gain ont tourné sur GPU, les trois qui n'en donnent pas sur CPU. Cinq essais ne
séparent pas ces deux causes (\cref{sec:llm-materiel}). Les trois premiers essais ont tourné avant
les garde-fous contre un calcul corrompu. Un gradient corrompu mais fini, écrêté à 1, a pu y passer
sans trace. Le calcul MPS n'est pas déterministe d'une tentative à l'autre, si bien que la même
graine aléatoire ne garantit pas le même essai sur GPU.

\section{Validité externe}
\label{sec:val-externe}

\subsection{La population des articles}

\begin{itemize}
  \item \textbf{La population a changé.} 27 articles du premier tirage de 2026 sur 50 traitent de
        grands modèles de langue, une famille absente en 2018. Nos verdicts valent pour le RL
        classique.
  \item \textbf{Le filtre de titre.} L'expression régulière appliquée aux titres rate les articles
        de RL dont le titre ne le dit pas, en 2018 comme en 2026. Mesuré pour 2026 avec un plan écrit
        d'avance (écart 24) : sur 400 articles tirés hors du filtre, 25 ont un résumé qui parle de
        RL, soit environ 1\,037 articles manqués (de 708 à 1\,503), 45~\% de l'ensemble. C'est un
        majorant, puisque la règle mécanique compte aussi des faux positifs. Lus avec les deux
        grilles, 20 de ces articles ressemblent au premier tirage : le code en lien dans 60~\% d'entre
        eux (67~\% dans le tirage), une incertitude dans 50~\% (47~\%). Les 7 articles de RL classique
        parmi eux sont trop peu nombreux pour éprouver nos verdicts.
  \item \textbf{Le tri sur le titre de la file de 2026.} Un article de RL classique dont le titre
        contient \enquote{reasoning}, \enquote{preference} ou \enquote{math} n'a pas pu entrer dans
        l'échantillon classique de 2026. Les 9 articles d'ICLR 2026 en tête de la file, injoignables
        le 9 novembre 2026, ont été sautés.
  \item \textbf{La source des PDF.} Les articles de 2026 ont été lus surtout dans leur version
        arXiv (45 articles de RL classique sur 50), ceux de 2018 dans les actes des conférences. Une
        version arXiv est parfois plus longue ou plus récente. Avec 5 articles seulement lus sur le
        site d'une conférence, cet effet ne se mesure pas.
  \item \textbf{La trajectoire sans ICLR.} Les années 2019 à 2024 ne couvrent qu'ICML et NeurIPS :
        OpenReview refuse les téléchargements automatisés.
  \item \textbf{Le tirage de Pineau.} Sa méthode de tirage de 2018 n'est pas décrite. Nos
        comparaisons avec ses chiffres supposent un tirage comparable au nôtre et servent de
        contexte, jamais de verdict.
\end{itemize}

\subsection{Les simulations}

Les expériences se limitent au simulateur MuJoCo et à quatre tâches de locomotion
(\task{HalfCheetah}, \task{Hopper}, \task{Walker2d}, \task{Swimmer}), celles de
\textcite{henderson2018matters}. Elles ne disent rien des jeux à images, de la robotique réelle ni du
RL hors ligne. E1 et B1 se limitent à une seule tâche (\task{HalfCheetah}) et à deux algorithmes
(PPO et TRPO). Les quatre programmes sont pris avec leurs hyperparamètres par défaut, sauf dans E6 :
un autre réglage pourrait changer les classements. La physique modifiée ne change que deux
grandeurs, les masses et le frottement.

\subsection{Les budgets}

PPO s'entraîne 1\,024\,000 pas, SAC et CrossQ 200\,000 (écarts 1 et 14 (d) des simulations). La lecture
à budget égal (écart 12 (iii)) compare PPO à 200\,000 pas, lu sur sa courbe d'un million de pas : un
PPO réglé pour un budget de 200\,000 pas pourrait faire mieux. Le budget de SAC et de CrossQ, réduit
pour tenir le calcul, les place dans le régime \enquote{peu de données}. Un budget plus long pourrait
changer leur ordre.

\subsection{Une seule plateforme}

Tous les essais viennent d'un seul Mac à puce Apple M2 Max, partagé avec d'autres projets, sous
macOS et pour arm64. MuJoCo, JAX et PyTorch arrondissent différemment d'un processeur à l'autre :
la même graine aléatoire donne une autre courbe sur une autre machine (\cref{sec:reentrainer}).

\subsection{L'échelle du volet GRPO}

Le volet GRPO reste une illustration : deux modèles de 0,5 et 1 milliard de paramètres, une tâche de
multiplication, 100 pas, des groupes de 8 réponses, des réponses limitées à 160 jetons, aucun terme
KL. \textcite{shao2025spurious} travaillent sur des modèles de 7 milliards de paramètres spécialisés
en mathématiques, sur MATH-500, avec bien plus de calcul. Notre volet ne réplique pas leur
expérience et ne la réfute pas.

\section{Validité de construit}
\label{sec:val-construit}

\begin{itemize}
  \item \textbf{Notre grille lit 2018 autrement que Pineau.} Notre lecture trouve 55~\% d'articles
        de 2018 qui donnent tous les hyperparamètres de leur méthode, quand Pineau en comptait
        90~\%. Les critères ont le même intitulé mais des règles de décision implicites différentes. C'est la
        raison de la règle de verdict à grille égale (écart 20).
  \item \textbf{Un critère satisfait n'est pas une pratique bien faite.} Q5 vaut oui pour un lien
        vers du code, sans vérifier que ce code tourne. X2 vaut oui pour une barre d'erreur, sans
        juger si elle mesure la bonne dispersion : la barre peut résumer des épisodes au lieu des
        graines aléatoires, ce que l'expérience B1 montre dangereux. Q8 vaut oui pour un test
        appliqué, sans juger s'il est approprié.
  \item \textbf{La mesure de performance (Y1)} est notre critère le plus fragile : 60~\% d'accord sur
        le premier tirage, avant que ses règles soient précisées. Le verdict de la faille F8 se lit
        avec cette réserve.
  \item \textbf{Les cas limites de X3, X4 et X7} n'ont reçu aucune règle d'harmonisation nouvelle,
        pour lire les deux années avec la même grille. Un calcul réduit au temps d'inférence ou des
        graines aléatoires ratées écartées restent codés comme le lecteur les a lus.
  \item \textbf{Ce qu'une liste de contrôle déclare.} Les réponses des auteurs à la liste de NeurIPS
        ne garantissent pas la pratique qu'elles nomment : les 18 articles qui répondent oui à la
        question sur la signification statistique n'appliquent aucun test. Le sondage mesure ce que
        l'article fait, la liste de contrôle ce qu'il déclare.
  \item \textbf{Le score d'un essai de simulation} est le retour moyen de dix épisodes déterministes
        au dernier pas. L'expérience B3 montre qu'une autre mesure (meilleur score, moyenne de la
        courbe) change l'ordre de certaines paires.
  \item \textbf{Le pass@1 du volet GRPO} mêle justesse et concision, puisqu'une réponse coupée
        compte pour fausse (\cref{sec:llm-coupure}).
  \item \textbf{Les notes de coût et de risque} qui ont mis à l'épreuve la thèse du coût viennent de
        trois modèles d'IA agentique, de trois tailles différentes. Leur accord entre eux ne
        garantit pas qu'un chercheur noterait de même.
\end{itemize}

\section{Validité de la conclusion statistique}
\label{sec:val-stat}

\paragraph{La taille des échantillons du sondage} Avec 50 articles de chaque côté, un changement de
dix points passe encore souvent inaperçu : \enquote{pas de changement détectable} ne veut pas dire
\enquote{aucun changement}. Le premier tirage comptait 14 articles de RL classique, trop peu pour
détecter un écart de trente points, d'où l'élargissement de l'écart 12. La trajectoire repose sur
20 articles par année : une part y a une incertitude d'environ vingt points de part et d'autre.

\paragraph{Les comparaisons multiples} Le sondage compare dix-sept critères dans la comparaison
principale. La correction de Holm protège le taux d'erreur de la famille, au prix de la puissance :
un écart dont l'intervalle de Newcombe évite zéro sans passer Holm, comme Q6 ou Y12, est rapporté
mais ne fait pas un verdict. Pour E2, 24 intervalles de probabilité d'amélioration sont aussi donnés
au niveau simultané de 95~\% par la correction de Bonferroni (écart 18 des simulations,
\cref{tab:e2-bonferroni}).

\paragraph{Les verdicts fragiles} La hausse de l'incertitude rapportée (X2) ne passe plus la
correction de Holm dans deux contrôles ($p = 0{,}07$ chaque fois). Elle reste le plus fragile des
quatre changements. Le code et le calcul passent le seuil dans tous nos contrôles, la sensibilité à
un hyperparamètre dans presque tous (\cref{tab:robustesse}).

\paragraph{Le nombre de graines aléatoires des simulations} Les résultats principaux reposent sur dix graines
aléatoires par couple. L'expérience E5 montre ce que ce nombre permet : entre deux programmes dont
les moyennes diffèrent de moins de 20~\%, cinq graines aléatoires ne retrouvent le bon ordre que 71 fois sur
cent en médiane, 62 fois mesuré sur vingt graines aléatoires. Les classements à dix graines sur les tâches où
le premier ne bat pas nettement le deuxième ne valent pas plus. Les vingt graines aléatoires de l'écart 12
confirment ou nuancent chaque résultat. Le \cref{chap:06-experiences-resultats} dit chaque
divergence.

\paragraph{Le nombre de graines aléatoires du volet GRPO} Avec cinq graines aléatoires, le bootstrap n'a que 126
rééchantillons distincts, dix avec trois graines aléatoires. Les intervalles du \cref{chap:07-modeles-de-langue}
sont très grossiers. Un intervalle qui contient zéro dit que nous ne pouvons pas conclure, pas que
l'effet est nul.

\paragraph{Les intervalles bootstrap} Les intervalles de rliable dépendent du générateur global de
NumPy et de l'ordre des appels. Leurs bornes bougent avec la graine aléatoire du bootstrap, de l'ordre du bruit
de Monte-Carlo (écart 13 des simulations) : un verdict dont la borne frôle le seuil reste fragile.

\section{Ce que testerait un nouvel entraînement sur une autre machine}
\label{sec:reentrainer}

Un nouvel entraînement complet sur un autre processeur ne redonnerait pas nos courbes. MuJoCo, JAX
et PyTorch arrondissent différemment d'un processeur à l'autre, si bien que la même graine aléatoire donne une
autre trajectoire dès les premiers pas. Le volet GRPO l'illustre : pour un même modèle et une même
graine aléatoire, le pass@1 de départ diffère entre GPU et CPU (\cref{sec:llm-materiel}).

Un tel entraînement testerait donc les statistiques, pas chaque courbe. Il dirait si la dispersion
entre graines aléatoires, les classements, les gonflements dus à la sélection et les pertes sous physique
modifiée restent du même ordre avec d'autres tirages. Une conclusion qui ne
survivrait qu'à nos propres graines aléatoires ne dirait rien des algorithmes. L'intégration continue fait déjà le test
équivalent pour les analyses : sur Linux, les comptes, les verdicts et les classements reviennent
exactement, les flottants à $10^{-12}$ près en relatif (\cref{sec:ci}). Aucun entraînement n'a été
refait sur une autre machine.
